% !TEX TS-program = pdflatexmk
\documentclass[12pt]{amsart}
\usepackage{multicol}
\usepackage{float}

%\usepackage[parfill]{parskip}    % Activate to begin paragraphs with an empty line rather than an indent

\usepackage[margin=1in]{geometry}


\usepackage{amsmath,amssymb,amsthm,latexsym,graphicx}
\usepackage[normalem]{ulem}
\usepackage{setspace} %used for doublespacing, etc.
\usepackage{hyperref}
\usepackage{cancel}
\usepackage[dvipsnames,usenames]{color}
\usepackage[all]{xy}
\usepackage{fancyhdr}
\pagestyle{fancy}
	\renewcommand{\headrulewidth}{0.5pt} % and the line
	\headsep=1cm
	
\DeclareGraphicsRule{.tif}{png}{.png}{`convert #1 `dirname #1`/`basename #1 .tif`.png}

%Some useful environments.
\newtheorem{theorem}{Theorem}
\newtheorem{corollary}[theorem]{Corollary}
\newtheorem{conjecture}[theorem]{Conjecture}
\newtheorem{lemma}[theorem]{Lemma}
\newtheorem{proposition}[theorem]{Proposition}
\newtheorem{definition}[theorem]{Definition}
\newtheorem{example}[theorem]{Example}
\newtheorem{axiom}{Axiom}
\theoremstyle{remark}
\newtheorem{remark}{Remark}
\newtheorem*{exercise}{Exercise}%[section]

%Some useful shortcuts for our favorite fields
\def\RR{\ensuremath{\mathbb R}} %Note, you can use these WITHOUT entering math mode
\def\NN{\ensuremath{\mathbb N}}
\def\ZZ{\ensuremath{\mathbb Z}}
\def\QQ{{\ensuremath\mathbb Q}}
\def\CC{\ensuremath{\mathbb C}}
\def\EE{{\ensuremath\mathbb E}}

%Some useful shortcuts for formatting lists
\newcommand{\bc}{\begin{center}}
\newcommand{\ec}{\end{center}}
\newcommand{\be}{\begin{enumerate}}
\newcommand{\ee}{\end{enumerate}}
\newcommand{\bi}{\begin{itemize}}
\newcommand{\ei}{\end{itemize}}

%Some useful shortcuts for formatting mathematical symbols
\newcommand{\ol}[1]{\overline{#1}}
\newcommand{\oimp}[1]{\overset{#1}{\iff}} %labeled iff symbol
\newcommand{\bv}[1]{\ensuremath{ \vec{\mathbf{#1}}} } %makes a vector.
\newcommand{\mc}[1]{\ensuremath{\mathcal{#1}}} %put something in caligraphic font
\newcommand{\mpg}[1]{\marginpar{ #1}} %to put comments in margins
\newcommand{\bsl}[1]{\texttt{\symbol{92}{\em #1}}} %for backslashes.
\newcommand{\normale}{\trianglelefteq}
\newcommand{\normal}{\triangleleft}

%Code for formatting the proofs a little nicer for submitted homework
\makeatletter
\renewenvironment{proof}[1][\proofname]{\par\doublespacing
  \pushQED{\qed}%
  \normalfont \topsep6\p@\@plus6\p@\relax
  \list{}{%
    \settowidth{\leftmargin}{\itshape\proofname:\hskip\labelsep}%
    \setlength{\labelwidth}{0pt}%
    \setlength{\itemindent}{-\leftmargin}%
  }%
  \item[\hskip\labelsep\itshape#1\@addpunct{:}]\ignorespaces
}{%
  \popQED\endlist\@endpefalse
  \singlespacing
}
\makeatother


%Commenting tools. You can ignore this stuff unless we're exchanging materials where I'm commenting in detail on how you are writing/using LaTeX.
\usepackage{soul}
\definecolor{highlight}{rgb}{1,0.6,0.6}
\sethlcolor{highlight}
\newcommand{\hlm}[1]{\colorbox{highlight}{$\displaystyle #1$}}
\newtheoremstyle{mycomment}{\topsep}{-0in}{\small \itshape \sffamily}{}{\small \itshape\sffamily}{:}{.5em}{}
\theoremstyle{mycomment}
\newtheorem*{acomment}{\color{BrickRed}{Comment}}
\newcommand{\com}[1]{{\color{OliveGreen}\begin{acomment}{#1} %#2 \color{black} 
\end{acomment}\noindent}}
%\newcommand{\com}[1]{{\color{BrickRed}{\\ Comment:}\color{OliveGreen}{#1} \\}}
\newcommand{\red}[1]{{\color{BrickRed} #1}}
\newcommand{\blue}[1]{{\color{MidnightBlue}#1}}
\newcommand{\green}[1]{{\color{OliveGreen}#1}}
\newcommand{\mwrong}[2]{\red{\cancel{#1}}\green{#2}}
\newcommand{\wrong}[2]{\red{\sout{#1}}\green{#2}}
\definecolor{OliveGreen}{rgb}{.3,.5,.2}
\definecolor{MidnightBlue}{rgb}{.3,.4,.6}
\newcommand{\pts}[1]{\hfill\blue{{#1}/5}}

\chead{MATH 631F}
\pagestyle{fancy}
%Modify these items:
\rhead{\emph{Algebra Group}}
\lhead{\emph{HW \#4 --- \today}}

\DeclareMathOperator{\lcm}{lcm}
\begin{document}

\thispagestyle{fancy}
\author{Coordinator: Robert Babac}


%\usepackage{float} <- Throw this in the preamble






\begin{exercise}[\textbf{3.1.41}] Let $G$ be a group. Prove that $N = \langle x^{-1}y^{-1}xy | x, y, \in G \rangle$ is a normal subgroup of $G$.
  and $G/N$ is abelian.\\
  \begin{proof} Let $G$ be a group, and suppose that $N = \langle x^{-1}y^{-1}xy | x, y, \in G \rangle$. Consider $g \in G$. Note that for the generator of $N$, 
    $gx^{-1}y^{-1}xyg^{-1} =  gx^{-1}(g^{-1}g)y^{-1}(g^{-1}g)x(g^{-1}g)yg^{-1} = (gx^{-1}g^{-1})(gy^{-1}g^{-1})(gxg^{-1})(gyg^{-1}) = (gxg^{-1})^{-1}(gyg^{-1})^{-1}(gxg^{-1})(gyg^{-1})$
    Let $gxg^{-1} = a$ and $gyg^{-1} = b$. Note that $gx^{-1}y^{-1}xyg^{-1} = a^{-1}b^{-1}ab$ therefore by definition $gx^{-1}y^{-1}xyg^{-1} \in N$. Thus $N \trianglelefteq G$.\\ 
    Since $x^{-1}y^{-1}xy \in N$ for all $x, y \in G$, and $N \trianglelefteq G$ by the previous exercise $G/N$ must be abelian. 
  \end{proof}

\end{exercise}
\vspace{.15in}




\begin{exercise}[\textbf{3.2.5}] Let $H$ be a subgroup of $G$ and fix some element $g \in G$.
  \begin{enumerate}
    \item[a] Prove that $gHg^{-1}$ is a subgroup of $G$ of the same order as $H$.\\
    \begin{proof} Suppose $H$ is a subgroup of $G$ and fix some element $g \in G$. Consider the group $gHg^{-1}$. Clearly this 
      group is non empty, as it contains the identity. Let $a, b \in gHg^{-1}$ for some $h_a, h_b \in H$ such that $a = gh_ag^{-1}$ and 
      $b = gh_bg^{-1}$. Note that $ab^{-1} = gh_ag^{-1}gh_b^{-1}g^{-1} = gh_ah_b^{-1}g^{-1}$. Since $H \leq G$ we know that $h_ah_b^{-1} \in H$ thus 
      $ab^{-1} \in gHg^{-1}$ and $gHg^{-1} \leq G$.\\
      To show that $|gHg^{-1}| = |H|$ we will prove they are isomorphic. Consider the map $\varphi: H \to gHg^{-1}$ defined by $\varphi(h) = ghg^{-1}$.
      Note that this map is a homomorphism, for all $a, b \in H$ we know that $\varphi(ab) = gabg^{-1} = gag^{-1}gbg^{-1} = \varphi(a)\varphi(b)$. This map is also an injection, if $\varphi(a) = \varphi(b)$ then we know that $gag^{-1} = gbg^{-1}$, and thus $a = b$. The map is a surjection by definition. Thus $gHg^{-1} \cong H$ and therefore $|gHg^{-1}| = |H|$.
    \end{proof}
    \vspace{.15in}
    \item[b] Deduce that if $n \in \ZZ^+$ and $H$ is the unique subgroup of $G$ of order $n$ then $H \trianglelefteq G$.\\
    \begin{proof} Suppose that $H$ is the unique subgroup of $G$ of order $n \in \ZZ^+$. Note that by the previous exercise the 
      subgroup $gHg^{-1}$ must have the same order of $H$, since $H$ has unique order we know $gHg^{-1} = H$. Thus  $H \trianglelefteq G$.
    \end{proof} 
  \end{enumerate}
\end{exercise}
\vspace{.15in}








  \end{document}